\documentclass[12pt,a4paper]{article}

% ---- Packages ----
\usepackage[margin=1in]{geometry}
\usepackage{amsmath}
\usepackage{amssymb}
\usepackage{array}
\usepackage{booktabs}
\usepackage{url}
\usepackage{hyperref}
\usepackage{graphicx}
\usepackage{setspace}
\usepackage{titlesec}
\usepackage{parskip}
\usepackage{microtype}
\usepackage[natbib=true, style=ieee, sorting=none]{biblatex}
\addbibresource{nanodb_references.bib}

% ---- Formatting ----
\onehalfspacing
\setlength{\parskip}{6pt}
\setlength{\parindent}{0pt}
\titleformat{\section}{\large\bfseries}{}{0em}{}[\titlerule]
\titleformat{\subsection}{\normalsize\bfseries}{}{0em}{}

\newcommand{\NanoDB}{\textsc{NanoDB}}
\newcommand{\BigO}[1]{$O(#1)$}
\newcommand{\code}[1]{\texttt{#1}}

\hypersetup{
    colorlinks=true,
    linkcolor=blue,
    citecolor=blue,
    urlcolor=blue,
    pdftitle={NanoDB: Design and Implementation of a Custom Relational Database Engine},
    pdfauthor={Qamar Mehmood}
}

% ======================================================
\begin{document}
% ======================================================

% ---- Title Block ----
\begin{center}
    {\LARGE \textbf{\NanoDB: Design and Implementation of a Custom}}\\[4pt]
    {\LARGE \textbf{Relational Database Engine with Constraint-Driven}}\\[4pt]
    {\LARGE \textbf{Data Structures}}\\[16pt]
    {\large Qamar Mehmood}\\
    {\normalsize Roll Number: 26I-7800}\\[4pt]
    {\normalsize CS-4002 Applied Programming \quad $\cdot$ \quad MS-CS Spring 2026}\\
    {\normalsize National University of Computer and Emerging Sciences (FAST-NUCES), Islamabad}\\
    {\normalsize Instructor: Bushra Fatima}\\[4pt]
    {\small GitHub: \url{https://github.com/Qamar2001/nanodb}}
\end{center}

\vspace{8pt}
\hrule
\vspace{16pt}

% ======================================================
\begin{abstract}
\noindent
Modern relational database management systems such as PostgreSQL and SQLite abstract their internals behind layers of library code, making them practical for production but opaque for learning. This paper presents the design, implementation, and empirical evaluation of \NanoDB, a fully functional mini-RDBMS built from scratch in C++ under a strict prohibition on the C++ Standard Template Library (STL). Every subsystem, from the page-based buffer pool to the query optimizer, was implemented using raw pointers and arrays. The engine loads 100,000 real-world records from the TPC-H decision-support benchmark \cite{tpch2023} across three relational tables, evaluates nested SQL expressions through a Shunting-Yard parser \cite{dijkstra1961} operating on a custom stack, and routes multi-table joins through a Kruskal MST cost optimizer \cite{kruskal1956}. Empirical measurements confirm a 69-fold query acceleration when using the AVL index \cite{avl1962} over a sequential scan at 20,000 records, and the buffer pool successfully managed 4,950 LRU eviction cycles \cite{oneil1993} under a constrained 50-page memory budget without memory corruption. This paper presents the architectural rationale, formal complexity proofs, and measured performance characteristics of each subsystem.
\end{abstract}

\vspace{8pt}
\hrule
\vspace{12pt}

% ======================================================
\section{Introduction}

The gap between \textit{using} a database and \textit{understanding} one is wide. A graduate student who writes complex JOIN queries every day may never stop to consider how the database decides which rows to examine, how a WHERE clause becomes a sequence of operations on stack-allocated temporaries, or why some join orderings are an order of magnitude more expensive than others. The \NanoDB{} project was designed to close that gap.

The defining constraint of this project is a complete ban on the C++ Standard Template Library. There is no \code{std::vector}, no \code{std::map}, no \code{std::stack}, and no \code{std::sort}. Every data structure must be re-engineered from first principles using raw arrays and pointer arithmetic. While this restriction may seem artificial, it forces a genuine engagement with algorithmic trade-offs that would otherwise remain hidden behind convenient interfaces. When you implement your own doubly-linked list and experience firsthand why \BigO{1} pointer relinking outperforms \BigO{N} array shifting by thousands of microseconds at realistic data scales, the lesson stays with you.

The engine that emerged from this process is a working relational system. It manages fixed-size memory pages through a custom buffer pool, indexes rows using a height-balanced AVL tree, parses SQL-like expressions using the classical Shunting-Yard algorithm on a hand-written stack, routes multi-table joins through a graph-based cost optimizer, serializes data to disk through raw binary I/O, and schedules high-priority administrative queries ahead of background reads through a binary max-heap. The dataset is 100,000 records from the TPC-H benchmark, the same industry standard used to evaluate commercial database systems \cite{tpch2023}.

This paper is organized as follows. Section 2 surveys the related academic work that underpins the design. Section 3 describes the overall system architecture. Sections 4 through 8 present each subsystem in detail with formal complexity analysis. Section 9 presents empirical benchmark results. Section 10 concludes with a reflection on what the implementation revealed.

% ======================================================
\section{Related Work}

The subsystems implemented in \NanoDB{} draw from several decades of foundational research in systems and algorithms.

\subsection{Buffer Management}
The idea of managing a fixed-size page pool in main memory and evicting pages to disk as needed is central to database systems design, as described extensively in Ramakrishnan and Gehrke \cite{ramakrishnan2003}. The Least Recently Used (LRU) replacement policy, used here, was analyzed in depth by O'Neil et al.\ \cite{oneil1993} in the context of their LRU-K algorithm, which extends basic LRU by tracking the $K$-th most recent access rather than just the most recent. \NanoDB{} implements the simpler LRU-1 variant, which is the standard textbook policy.

\subsection{Balanced Tree Indexing}
The AVL tree was introduced by Adelson-Velsky and Landis in 1962 \cite{avl1962}. It was the first self-balancing binary search tree and remains one of the most pedagogically clear, with an explicit height-difference invariant of at most one between any node's left and right subtrees. Alternative structures such as red-black trees (used in \code{std::map}) and B+ trees (used in most commercial databases) trade some simplicity for better cache behavior. For this project, the AVL tree's clear invariant and well-understood rotation cases make it a natural choice. The general treatment of balanced BSTs and their complexity proofs appears in Cormen et al.\ \cite{cormen2009}.

\subsection{Expression Parsing}
The Shunting-Yard algorithm was described by Dijkstra in 1961 \cite{dijkstra1961} as a method for converting mathematical expressions from infix to postfix (Reverse Polish Notation) notation. The resulting postfix stream can be evaluated in a single linear pass using a stack, with no backtracking or lookahead. The algorithm handles operator precedence and left-to-right associativity through simple comparisons of precedence values on the operator stack. \NanoDB{} extends the algorithm to handle logical operators AND and OR alongside arithmetic and comparison operators.

\subsection{Query Optimization}
The field of cost-based query optimization was formalized by Selinger et al.\ \cite{selinger1979} in the IBM System R optimizer, which introduced the concept of generating all possible join orderings and selecting the one with the lowest estimated cost. Full System R-style dynamic programming is complex to implement. \NanoDB{} uses a simplified model: join orderings are represented as edges in a weighted graph and Kruskal's minimum spanning tree algorithm \cite{kruskal1956} selects the cheapest non-cyclic set of joins. The Union-Find data structure used to detect cycles during Kruskal's execution was analyzed for efficiency by Tarjan \cite{tarjan1975}.

% ======================================================
\section{System Architecture}

\NanoDB{} is organized as a layered pipeline where each component depends on the one below and exposes a clean interface to the one above.

At the lowest level sits the \textbf{Pager}, which maps table data to fixed-size binary pages on disk. The pager delegates in-memory residency tracking to the \textbf{LRU Cache}, which uses a doubly-linked list combined with a hash map to achieve constant-time page lookup, promotion, and eviction.

Above the storage layer sits the \textbf{Schema Engine}. Rows are heterogeneous: each row owns a raw \code{DataType**} where each cell points to an \code{IntType}, \code{FloatType}, or \code{StringType} derived from a common abstract base class. Virtual comparison operators and a virtual \code{clone()} method allow the rest of the engine to manipulate row data without any type-checking if-statements.

The \textbf{Hash Map Catalog} sits between the schema and the executor. It maps table names to their corresponding \code{TableEntry} objects using separate chaining with the djb2 hash function, rehashing automatically when the load factor exceeds 0.75.

The \textbf{Query Parser} accepts a raw SQL-like string, tokenizes it, extracts the WHERE clause, and converts it to a postfix token stream using the Shunting-Yard algorithm on a custom \code{Stack<int>}. The postfix stream is stored in a \code{TokenStream} object and later evaluated by the executor.

The \textbf{Executor} connects everything. For a SELECT query, it resolves table names through the catalog, uses the AVL primary-key index for point lookups or falls back to a sequential scan for range queries, and evaluates the postfix WHERE expression against each candidate row. For multi-table queries, it builds a weighted graph over the involved tables and computes a minimum spanning tree to determine the cheapest join ordering before executing a nested-loop join.

Finally, a \textbf{Priority Queue} wraps the execution pipeline. Queries tagged as \code{ADMIN} receive priority value 100 and are dequeued before standard queries with priority 0, regardless of submission order.

% ======================================================
\section{Memory Management: Buffer Pool and LRU Cache}

\subsection{Design Rationale}
A database cannot load an arbitrarily large table entirely into RAM \cite{ramakrishnan2003}. The buffer pool manages a fixed-size set of memory pages, loading them from disk on demand and evicting the least recently used page when the pool is full. Dirty pages, those modified since being loaded, must be written back to disk before eviction. Clean pages can be discarded directly.

\subsection{Implementation}
The LRU cache is built on two hand-coded components working together. A doubly-linked list maintains usage ordering: the head holds the most recently accessed page, and the tail holds the least recently accessed. When a page is accessed, it is moved to the head in \BigO{1} time by relinking its predecessor and successor pointers. When eviction is required, the tail is removed and its page is flushed to disk if dirty.

A hash map (described in Section 6) maps page identifiers to their list nodes, allowing the cache to determine whether a page is resident in \BigO{1} expected time without traversing the list.

\subsection{Formal Complexity Analysis}
Let $C$ denote the cache capacity and $B$ the number of hash buckets.

\textbf{Space:} $O(C + B)$. The list holds at most $C$ nodes and the hash map holds at most $C$ entries across $B$ chains.

\textbf{Lookup:} $O(1)$ expected. The djb2 hash function distributes keys uniformly, giving expected chain length $C/B$, a small constant when $B$ is proportional to $C$.

\textbf{Eviction:} $O(1)$ strict. Relinking the tail requires changing exactly two pointers regardless of list size.

\textbf{Promotion:} $O(1)$ strict. Moving a node to the head requires relinking its predecessor, its successor, and the existing head node. No traversal is required because the hash map provides a direct pointer to the node.

These guarantees mean the buffer pool adds zero asymptotic overhead to query processing regardless of dataset scale \cite{oneil1993}.

% ======================================================
\section{Indexing: AVL Self-Balancing Tree}

\subsection{The Cost of Linear Scans}
A sequential scan answering \code{WHERE c\_custkey == 10000} on the 20,000-row customer table must examine up to 20,000 rows in the worst case. The AVL tree index answers the same query in at most $\lceil 1.44 \log_2(20002) \rceil = 21$ comparisons \cite{avl1962}.

\subsection{Implementation}
Each AVL tree node stores an integer key, a row identifier, a height attribute, and left and right child pointers. Four rotation cases restore the height-balance invariant after insertions that violate it: LL (single right rotation), RR (single left rotation), LR (left-then-right rotation), and RL (right-then-left rotation). Each rotation involves a fixed number of pointer assignments and one height update, taking strictly $O(1)$ time.

\subsection{Formal Complexity Analysis}
Let $N_h$ denote the minimum number of nodes in an AVL tree of height $h$. The recurrence $N_h = 1 + N_{h-1} + N_{h-2}$ with base cases $N_0 = 1$, $N_1 = 2$ grows at least as fast as the Fibonacci sequence. By Fibonacci's closed form, $h < 1.44 \log_2(N + 2)$ for a tree with $N$ nodes \cite{cormen2009}.

\textbf{Space:} $\Theta(N)$, one node per indexed record.

\textbf{Search:} $O(\log N)$. Each comparison at a node eliminates one subtree and the bounded height guarantees termination within $1.44 \log_2(N+2)$ steps.

\textbf{Insertion:} $O(\log N)$. One root-to-leaf path of length $O(\log N)$ followed by at most one rotation.

% ======================================================
\section{Hash Map System Catalog}

The system catalog maps table names (strings) to \code{TableEntry} pointers. Because the executor resolves table names on every query execution, \BigO{N} linear search is unacceptable even for small catalogs.

\NanoDB{} implements a chaining hash map using the djb2 polynomial rolling hash \cite{ramakrishnan2003}:
\[
  h = h \times 33 + c_i
\]
where $c_i$ is the ASCII value of the $i$-th character. djb2 produces good distribution across table name inputs. The map resizes (rehashes) automatically when the load factor $\alpha = N/B$ exceeds 0.75, keeping expected chain lengths bounded.

\textbf{Space:} $\Theta(B + N)$ for $B$ buckets and $N$ entries.
\textbf{Lookup/Insert/Delete:} $O(1)$ expected, $O(N)$ worst case (all keys in one chain).

% ======================================================
\section{Query Parsing and Expression Evaluation}

\subsection{Tokenization}
The tokenizer scans the WHERE clause character by character, producing a typed token stream: column identifiers, integer and floating-point literals, string literals in double quotes, arithmetic operators (\code{+}, \code{-}, \code{*}, \code{/}, \code{\%}), comparison operators (\code{==}, \code{!=}, \code{<}, \code{>}, \code{<=}, \code{>=}), logical operators \code{AND} and \code{OR}, and parentheses.

\subsection{Shunting-Yard Conversion}
The Shunting-Yard algorithm \cite{dijkstra1961} processes the infix token stream with two auxiliary structures: an output queue and a custom operator stack. Operand tokens go directly to the output. Operator tokens are held on the stack until an operator of lower or equal precedence must be pushed, at which point higher-precedence operators are popped to the output. Right parentheses trigger popping until the matching left parenthesis is consumed. The precedence table assigns: OR = 1, AND = 2, comparisons = 3, addition/subtraction = 4, multiplication/division/modulo = 5. The result is a postfix token stream stored in a \code{TokenStream} object.

\subsection{Postfix Evaluation}
At query time, the Executor evaluates the postfix stream against each candidate row using a custom \code{Stack<DataType*>}. Operand tokens are materialized into \code{DataType*} heap objects (looked up from the row for column references, or constructed fresh for literals). Binary operators pop two operands, compute the result, and push a new \code{DataType*} back. Because all types share virtual comparison and arithmetic operators, the evaluator handles integers, floats, and strings uniformly.

\subsection{Formal Complexity}
For $T$ tokens: tokenization is $\Theta(|$input$|)$ in characters; Shunting-Yard is $\Theta(T)$ since each token is pushed and popped at most once; postfix evaluation is $\Theta(T)$ for the same reason. Total parsing and evaluation is $O(T)$ \cite{dijkstra1961}.

% ======================================================
\section{Query Optimizer and Join Scheduling}

\subsection{The Join Ordering Problem}
For a three-table join customer $\bowtie$ orders $\bowtie$ lineitem, six execution orderings are possible. Joining the smaller tables first produces smaller intermediate results and reduces total work, but choosing the right order requires an estimate of intermediate result sizes \cite{selinger1979}.

\subsection{Graph-Based MST Optimizer}
\NanoDB{} models the join ordering problem as a minimum spanning tree problem. Each involved table is a graph node. Each pair of tables is connected by an edge whose weight is the product of their row counts, a proxy for the intermediate result size if those two tables are joined directly. Kruskal's algorithm \cite{kruskal1956} processes edges in ascending weight order, adding each edge to the MST if it does not form a cycle. Cycle detection uses a Union-Find structure with path compression and union by rank, which runs in near-constant inverse Ackermann time $\alpha(V)$ per operation \cite{tarjan1975}.

The MST join order is logged before execution, for example: \texttt{[LOG] Multi-table join routed via MST: customer -> orders -> lineitem}. This allows the evaluator to verify the optimizer decision at demo time.

\subsection{Priority Queue Scheduling}
A binary max-heap \code{PriorityQueue<T>} wraps the query pipeline. Each query object carries an integer priority: 100 for ADMIN queries and 0 for standard queries. The heap ensures that the highest-priority query is always at the root. An ADMIN query submitted after 50 background SELECTs will still be dequeued and executed first, with the event logged as: \texttt{[LOG] Admin query intercepted before background reads}.

\textbf{Heap Insert/Extract-Max:} $O(\log N)$ where $N$ is the number of pending queries \cite{cormen2009}.

% ======================================================
\section{Empirical Evaluation}

\subsection{Experimental Setup}
All benchmarks were compiled with \code{g++ -std=c++17 -O2} on Windows 11 (AMD processor). Timing used C++ \code{<chrono>} high-resolution clock at nanosecond granularity. The TPC-H dataset \cite{tpch2023} was loaded at full project scale: 20,000 customer rows, 30,000 orders rows, and 50,000 lineitem rows.

\subsection{Indexed vs. Sequential Scan}

The benchmark \code{benchmark\_runner.cpp} builds tables of 1,000, 10,000, and 100,000 rows and measures the time to locate the last key both by a sequential scan (worst case: the target is the final row) and by AVL index lookup.

\begin{table}[h]
\centering
\caption{Sequential scan vs.\ AVL index lookup times (nanoseconds, worst-case target key).}
\begin{tabular}{rrrr}
\toprule
\textbf{Records} & \textbf{Sequential Scan (ns)} & \textbf{AVL Index (ns)} & \textbf{Speedup} \\
\midrule
1,000    & 2,800     & 200   & $14.0\times$ \\
10,000   & 44,200    & 200   & $221.0\times$ \\
100,000  & 1,494,500 & 1,700 & $879.1\times$ \\
\bottomrule
\end{tabular}
\end{table}

The sequential scan grows linearly with record count, consistent with $O(N)$ behavior. The AVL index time grows logarithmically, consistent with $O(\log N)$: from 1,000 to 100,000 records (a $100\times$ size increase), AVL lookup time increases from 200 ns to only 1,700 ns ($8.5\times$). As predicted by $\log_2(100000) / \log_2(1000) \approx 5.56$, tree height grows slowly. The difference in speedup factors across scales confirms that the AVL advantage is not constant but increases as the dataset grows, matching the theoretical gap between $O(N)$ and $O(\log N)$ \cite{avl1962, cormen2009}.

The live demo result for a 20,000-row customer lookup was: sequential scan 131,900 ns, AVL index 1,900 ns, a $69.4\times$ speedup.

\subsection{LRU Memory Stress Test}

The buffer pool capacity was fixed at $P = 50$ pages and subjected to $M$ unique sequential page accesses. Because the access pattern is strictly increasing (no page is revisited), every access after the first 50 triggers an eviction. The expected eviction count is $\max(0, M - P)$.

\begin{table}[h]
\centering
\caption{LRU eviction counts under constrained 50-page buffer pool.}
\begin{tabular}{rrr}
\toprule
\textbf{Unique Pages Touched} & \textbf{LRU Evictions} & \textbf{Eviction Rate} \\
\midrule
1,000  & 950   & 95.0\% \\
5,000  & 4,950 & 99.0\% \\
10,000 & 9,950 & 99.5\% \\
\bottomrule
\end{tabular}
\end{table}

All 4,950 evictions in the project demo run completed without memory corruption, double-free errors, or data loss. Because each eviction runs in $O(1)$ time through the doubly-linked list \cite{oneil1993}, the total eviction overhead was negligible compared to disk I/O time.

\subsection{Demo Test Cases A through G}

\textbf{A: Parser and Evaluator.} The query \code{(c\_acctbal > 5000 AND c\_mktsegment == "BUILDING") OR c\_nationkey == 15} was successfully parsed into postfix form by the Shunting-Yard algorithm and evaluated row-by-row, with the postfix token sequence logged.

\textbf{B: Index Optimizer.} Both sequential scan and AVL lookup were applied to a specific \code{c\_custkey} in the 100,000-record dataset. Both timings were printed to the console and the speedup confirmed visually.

\textbf{C: Join Optimizer.} A three-table customer/orders/lineitem join was executed. The MST optimizer logged the chosen join path before the nested-loop join produced the correct matched rows.

\textbf{D: Memory Stress.} A 50-page buffer pool processed 5,000 distinct page accesses and printed 4,950 LRU evictions, exactly matching the formula $\max(0, 5000 - 50) = 4950$.

\textbf{E: Priority Queue.} Fifty standard SELECT queries were enqueued, followed by one ADMIN INSERT. The priority queue dequeued and executed the admin transaction first, confirmed by the execution log.

\textbf{F: Deep Expression Tree.} The query \code{((o\_totalprice * 1.5) > 100000 AND (o\_custkey \% 2 == 0)) OR (o\_orderstatus != "O")} exercised arithmetic, modulo, comparison, AND, OR, and nested parentheses simultaneously through postfix evaluation.

\textbf{G: Durability and Persistence.} Five customer records were inserted and all tables were serialized to binary files under \code{data/}. After a full engine restart with the \code{--reload} flag, the five records were successfully recovered and AVL indexes were rebuilt from the loaded rows.

% ======================================================
\section{Conclusion}

\NanoDB{} demonstrates that a fully functional relational database engine, with its own memory manager, type system, SQL parser, indexer, join optimizer, and persistence layer, can be built from first principles in C++ without any standard library containers. The empirical results confirm the theoretical predictions derived from algorithmic analysis: the AVL index delivers $879\times$ faster lookups than a sequential scan at 100,000 records, the LRU cache manages thousands of eviction cycles in constant time per eviction, and the MST optimizer correctly routes multi-table joins through the cheapest available join path.

Beyond the measurements themselves, the value of this project lies in what the implementation reveals. Design decisions that seem arbitrary in a high-level framework, such as why databases use doubly-linked lists rather than arrays for LRU, why AVL trees maintain a height invariant rather than simply sorting keys, or why query planners go to such lengths to compute optimal join orderings, become immediately obvious once you have experienced the cost of the wrong choice firsthand. The 4,950 eviction cycles that ran in negligible time would have cost millions of element shifts had an array been used instead. The 879-fold index speedup would disappear if the tree lost its height invariant. These lessons do not come from reading documentation; they come from implementing the system yourself.

% ======================================================
\printbibliography[title={References}]

\end{document}
